% ============================================================
% Ejercicio 1
% ============================================================
\section{Ejercicio 1}

\begin{tcolorbox}[enunciado]
\begin{enumerate}[label=(\alph*), leftmargin=*]
    \item ¿Cómo podemos expresar este problema en su versión de búsqueda?
    \item ¿Tiene sentido plantearse una versión de optimización?
    \item Plantea el problema en su versión de decisión, y enúncialo en su forma canónica.
\end{enumerate}
\end{tcolorbox}

\subsection{Solución}
\begin{enumerate}[label=(\alph*), leftmargin=*]
    \item \textbf{Un problema de búsqueda} consiste es encontrar y devolver la estructura exacta que cumple con las condiciones establecidas.

Para el caso de ``DeliCiencias'', la versión de búsqueda se expresaría formalmente de la siguiente manera:

Dado un catálogo de pedidos pendientes donde cada pedido $i$ tiene un peso físico $w_i$ y una ganancia económica $v_i$, una capacidad máxima de mochila $W$ y una ganancia meta $V$, \textbf{encuentra y devuelve} un subconjunto específico de pedidos $S$ tal que la suma de sus pesos no exceda $W$ ($\sum_{i \in S} w_i \le W$) y la suma de sus ganancias alcance o supere la meta $V$ ($\sum_{i \in S} v_i \ge V$). Si no existe ninguna combinación que cumpla ambas condiciones, el algoritmo debe reportarlo.

    \item Si, tiene sentido, ya que se busca maximizar la ganancia del repartidor aprovechando al máximo la capacidad $W$ de la mochila, sin depender de un umbral fijo $V$.
    
    La versión de optimización modificaría la estructura del problema al eliminar el parámetro de entrada $V$. El algoritmo ya no se limitaría a buscar una combinación que simplemente cruce un umbral, sino que evaluaría el catálogo para encontrar el subconjunto exacto de pedidos cuya suma de pesos no exceda $W$ y cuya suma de ganancias sea el máximo absoluto posible. 
    
    \item \textbf{OPTIMIZACIÓN DE ENTREGAS - VERSIÓN DE DESICIÓN.}\\
    \textbf{Datos (parámetros):} Un conjunto finito de \underline{pedidos pendientes} $P = \{p_1, p_2, \dots, p_n\}$. Para cada pedido $p_i \in P$, un \underline{peso físico} $w_i \in \mathbb{Z}^+$ y una \underline{ganancia económica} $v_i \in \mathbb{Z}^+$. Una \underline{capacidad máxima} de la mochila $W \in \mathbb{Z}^+$ y una \underline{ganancia meta} $V \in \mathbb{Z}^+$.
    
    \textbf{Pregunta (objetivo):} ¿Existe un subconjunto $S \subseteq P$ tal que la suma de los pesos de los pedidos en $S$ sea menor o igual a $W$ ($\sum_{p_i \in S} w_i \le W$) y la suma de sus ganancias sea mayor o igual a $V$ ($\sum_{p_i \in S} v_i \ge V$)?
\end{enumerate}